\scanpage{77}
\mspage{26}

Par contre, on s'attend que les \deleted{\unsure{notions}} morphismes $W$-asphériques soient stables par chgt. de base $W$-\uline{lisses} (qui généralisent les iso. locaux), et les morphismes $W$-coasphériques par chg. de base $W$-\unsure{propres|propre}. Cf plus bas.

\uline{Corollaire 2} Soit $f:X\longrightarrow Y$ \inserted{dans $\mathrm{Cat}$} une \uline{équivalence} de catégories. Alors $f$ est $W$-asphérique et $W$-coasphérique.

En effet, $f$ a un adjoint à \unsure{g} et un adjoint à droite, et on utilise la remarque page 23.

\uline{Prop. 3}. Soit $f:X\longrightarrow Y$ de $\mathrm{Cat}$. Conditions équivalentes :

a) Pour tout $y_0\in\mathrm{Ob}\,Y$, le foncteur d'inclusion \[X_{y_0}\mathrel{\relbar}X_{/y_0}\tag*{(\unsure{*|x})}\] est $W$-\uline{co}asphérique.

b) Pour tout $x_0\in X$, le morphisme induit \unsure{par les colocalisés|pour les colocalisés} \[{}_{x_0}\backslash X\mathrel{\relbar}{}_{y_0}\backslash Y\qquad(\text{où }y_0=f(x_0))\] est à fibres $W$-asphériques.

c) Pour tt \deleted{$\unsure{\Delta}$} morphisme $\mathord{\tikz[baseline=0pt,x=1ex,y=1ex]{\draw[line width=.55pt,line join=round] (0,0)--(.78,1.4)--(1.53,0)--cycle; \draw[line width=.55pt] (.25,0)--(.90,1.18);}}^1=Y'\longrightarrow Y$, considérant $X'=X\times_Y Y'$ sur $\mathord{\tikz[baseline=0pt,x=1ex,y=1ex]{\draw[line width=.55pt,line join=round] (0,0)--(.78,1.4)--(1.53,0)--cycle; \draw[line width=.55pt] (.25,0)--(.90,1.18);}}^1$, l'inclusion $X'_1\longrightarrow X'$ de la fibre fermée de $X'$ (fibre sur le point fermé de $Y'$) est $W$-coasphérique.

d) Pour toute flèche $u_0:y_0\longrightarrow y_1$ dans $Y$, et tout $x_0$ de $X_{y_0}$, la catégorie $\unsure{X|\mathscr X}(x_0,u_0)$ formée des flèches \inserted{$u:x_0\to x$} dans $X$ de source $x_0$, et qui relèvent $u_{\unsure{0}}$, est $W$-asphérique.

% archive sheet number: 77
